% Partie : logique du premier ordre
% Chapitre : déduction axiomatique
% Section : propriétés de la déductibilité
% Propriétés nécessaires aux ensembles maximalement cohérents.

\documentclass[../../../include/open-logic-section]{subfiles}

\begin{document}

\iftag{FOL}
      {\olfileid[fr]{fol}{axd}{pva}}
      {\olfileid[fr]{pl}{axd}{pva}}

\olsection{Propriétés de la \usetoken{s}{derivability}}

\begin{editorial}
Cette section présente une formulation plus ancienne du calcul
axiomatique. Ses désignations \textbf{Ax0} à \textbf{Ax3} renvoient
aux schémas suivants, ainsi qu'à leurs clôtures universelles :
\textbf{Ax0}, les instances des tautologies propositionnelles ;
\textbf{Ax1}, $\lforall[x][!B]\lif !B(t)$, lorsque $t$ est libre
pour $x$ dans $!B$ ;
\textbf{Ax2},
$\lforall[x][(!B\lif !C)]\lif
(\lforall[x][!B]\lif\lforall[x][!C])$ ;
\textbf{Ax3}, $!B\lif\lforall[x][!B]$, lorsque $x$ n'est pas
libre dans $!B$. La règle de construction des démonstrations est
le modus ponens. Si le quantificateur existentiel est primitif,
on suppose aussi le schéma dual
$!B(t)\lif\lexists[x][!B]$, sous la même condition de
substitution ; s'il est défini par la négation et le quantificateur
universel, ce schéma en découle. Ces conventions explicitent
les schémas invoqués
par la source ; elles ne sont pas les règles de quantification
présentées dans la formulation principale du chapitre.
\end{editorial}

\begin{prop}[Monotonie]
Si $\Gamma \subseteq \Delta$ et $\Gamma \Proves !A$, alors
$\Delta \Proves !A$.
\end{prop}

\begin{proof}
Tout sous-ensemble fini $\Gamma_0 \subseteq \Gamma$ est aussi un
sous-ensemble fini de~$\Delta$. Une !!{derivation} de
$\Gamma_0 \Proves !A$ établit donc également
$\Delta \Proves !A$.
\end{proof}

\begin{prop}\ollabel{prop:provability}
$\Gamma \Proves !A$ si et seulement si
$\Gamma\cup \{\lnot!A \}$ est incohérent.
\end{prop}

\begin{prop}
\begin{enumerate}
\item \ollabel{prop:provability-contr}
Si $\Gamma \Proves !A$ et
$\Gamma \cup \{!A\} \Proves \lfalse$, alors $\Gamma$ est incohérent.

\item \ollabel{prop:provability-lnot}
Si $\Gamma \cup \{!A\} \Proves \lfalse$, alors
$\Gamma \Proves \lnot !A$.

\item \ollabel{prop:provability-exhaustive}
Si $\Gamma \cup \{!A\} \Proves \lfalse$ et
$\Gamma \cup \{\lnot !A\} \Proves \lfalse$, alors
$\Gamma \Proves \lfalse$.

\tagitem{prvOr}{\ollabel{prop:provability-lor-left}
Si $\Gamma \cup \{!A\} \Proves \lfalse$ et
$\Gamma \cup \{!B\} \Proves \lfalse$, alors
$\Gamma \cup \{!A \lor !B\} \Proves \lfalse$.}{}

\tagitem{prvOr}{\ollabel{prop:provability-lor-right}
Si $\Gamma \Proves !A$ ou $\Gamma \Proves !B$, alors
$\Gamma \Proves !A \lor !B$.}{}

\tagitem{prvAnd}{\ollabel{prop:provability-land-left}
Si $\Gamma \Proves !A \land !B$, alors
$\Gamma \Proves !A$ et $\Gamma \Proves !B$.}{}

\tagitem{prvAnd}{\ollabel{prop:provability-land-right}
Si $\Gamma \Proves !A$ et $\Gamma \Proves !B$, alors
$\Gamma \Proves !A \land !B$.}{}

\tagitem{prvIf}{\ollabel{prop:provability-mp}
Si $\Gamma \Proves !A$ et $\Gamma \Proves !A \lif !B$,
alors $\Gamma \Proves !B$.}{}

\tagitem{prvIf}{\ollabel{prop:provability-lif}
Si $\Gamma \Proves \lnot !A$ ou $\Gamma \Proves !B$, alors
$\Gamma \Proves !A \lif !B$.}{}
\end{enumerate}
\end{prop}

\begin{proof}
Dans chaque cas, choisissons les sous-ensembles finis de~$\Gamma$
qui contiennent les hypothèses effectivement utilisées.
\begin{enumerate}
\item Supposons que $\Gamma_0 \cup \{!A\} \Proves \lfalse$
et que $\Gamma_1 \Proves !A$.

\begin{derivation}
1. & $\Gamma_0 \cup \{!A\} \Proves \lfalse$ & hyp. \\
2. & $\Gamma_1 \Proves !A$ & hyp. \\
3. & $\Gamma_0 \cup \Gamma_1 \Proves \lfalse$ & coupure \\
\end{derivation}

Comme $\Gamma_0 \subseteq \Gamma$ et $\Gamma_1 \subseteq \Gamma$,
on a $\Gamma_0 \cup \Gamma_1 \subseteq \Gamma$, d'où
$\Gamma \Proves \lfalse$.

\item Supposons que $\Gamma_0 \cup\{!A\} \Proves \lfalse$.\footnote{
Correction de l'édition française : la source suppose ici
la contradiction à partir de $\lnot !A$ et conclut $!A$,
alors que l'énoncé demande de partir de $!A$ et de conclure
$\lnot !A$. Les quatre lignes sont rétablies dans ce sens.}

\begin{derivation}
1. & $\Gamma_0 \cup\{!A\} \Proves \lfalse$ & hyp. \\
2. & $\Gamma_0 \Proves !A \lif \lfalse$ & théorème de la déduction \\
3. & $\Gamma_0 \Proves (!A \lif \lfalse) \lif \lnot !A$ & Ax0 \\
4. & $\Gamma_0 \Proves \lnot !A$ & MP 2, 3 \\
\end{derivation}

La monotonie donne $\Gamma \Proves \lnot !A$.

\item Supposons que $\Gamma_0 \cup \{!A\} \Proves \lfalse$
et que $\Gamma_1 \cup \{\lnot !A\} \Proves \lfalse$.

\begin{derivation}
1. & $\Gamma_0 \cup\{!A\} \Proves \lfalse$ & hyp. \\
2. & $\Gamma_0 \Proves !A \lif \lfalse$ & théorème de la déduction \\
3. & $\Gamma_1 \cup \{\lnot !A\} \Proves \lfalse$ & hyp. \\
4. & $\Gamma_0 \Proves (!A \lif \lfalse) \lif \lnot !A$ & Ax0 \\
5. & $\Gamma_0 \Proves \lnot !A$ & MP 2, 4 \\
6. & $\Gamma_0 \cup \Gamma_1 \Proves \lfalse$ & coupure 3, 5 \\
\end{derivation}

Comme $\Gamma_0 \subseteq \Gamma$ et $\Gamma_1 \subseteq \Gamma$,
on a $\Gamma_0 \cup \Gamma_1 \subseteq \Gamma$. Ainsi,
$\Gamma \Proves \lfalse$.

% prop:provability-lor-left
\tagitem{defOr}{}{
\iftag{probOr}{Exercice.}{Exercice.}}

% prop:provability-lor-right
\tagitem{defOr}{}{\iftag{probOr}{Exercice.}{
Supposons que $\Gamma_0 \Proves !A$.

\begin{derivation}
1. & $\Gamma_0 \Proves !A$ & hyp. \\
2. & $\Gamma_0 \Proves !A \lif (!A \lor !B)$ & Ax0 \\
3. & $\Gamma_0 \Proves !A \lor !B$ & MP 1, 2 \\
\end{derivation}

Donc $\Gamma \Proves !A \lor !B$. Le cas où
$\Gamma \Proves !B$ se traite de la même façon.}}

% prop:provability-land-left
\tagitem{defAnd}{}{\iftag{probAnd}{Exercice.}{
Supposons que $\Gamma_0 \Proves !A \land !B$.

\begin{derivation}
1. & $\Gamma_0 \Proves !A \land !B$ & hyp. \\
2. & $\Gamma_0 \Proves (!A\land !B) \lif !A$ & Ax0 \\
3. & $\Gamma_0 \Proves !A$ & MP 1, 2 \\
\end{derivation}

La monotonie donne $\Gamma \Proves !A$. Une !!{derivation}
analogue établit $\Gamma \Proves !B$.\footnote{Correction de
l'édition française : la troisième ligne de la source conclut
à tort $!A\lor !B$ au lieu de $!A$.}}}

% prop:provability-land-right
\tagitem{defAnd}{}{\iftag{probAnd}{Exercice.}{Exercice.}}

% prop:provability-mp
\tagitem{defIf}{}{\iftag{probIf}{Exercice.}{
Supposons que $\Gamma_0 \Proves !A$ et
$\Gamma_0 \Proves !A \lif !B$.

\begin{derivation}
1. & $\Gamma_0 \Proves !A$ & hyp. \\
2. & $\Gamma_0 \Proves !A \lif !B$ & hyp. \\
3. & $\Gamma_0 \Proves !B$ & MP 1, 2 \\
\end{derivation}

Comme $\Gamma_0 \subseteq \Gamma$, nous obtenons
$\Gamma \Proves !B$.}}

% prop:provability-lif
\tagitem{defIf}{}{\iftag{probIf}{Exercice.}{
Supposons d'abord que $\Gamma \Proves \lnot !A$.

\begin{derivation}
1. & $\Gamma_0 \Proves \lnot !A$ & hyp. \\
2. & $\Gamma_0 \Proves \lnot !A \lif (!A \lif !B)$ & Ax0 \\
3. & $\Gamma_0 \Proves !A \lif !B$ & MP 1, 2 \\
\end{derivation}

Le cas $\Gamma \Proves !B$ est analogue :

\begin{derivation}
1. & $\Gamma_0 \Proves !B$ & hyp. \\
2. & $\Gamma_0 \Proves !B \lif (!A \lif !B)$ & Ax0 \\
3. & $\Gamma_0 \Proves !A \lif !B$ & MP 1, 2 \\
\end{derivation}}}

\end{enumerate}
\end{proof}

\begin{probtag}{probOr,probAnd,probIf}
Achever la démonstration de la
\olref[fol][axd][pva]{prop:provability}.
\end{probtag}

\begin{thm}[Généralisation]\ollabel{thm:Generalization}
Si $\Gamma \Proves !A$ et si $x$ n'est libre dans aucune
!!{formula} de $\Gamma$, alors
$\Gamma \Proves \lforall[x][!A]$.
\end{thm}

\begin{proof}
Procédons par récurrence sur la construction de
$\mathsf{Thm}_\Gamma$, l'ensemble des formules déductibles
de~$\Gamma$. Si $!A$ est un axiome, sa clôture
$\lforall[x][!A]$ est également un axiome. Si $!A\in\Gamma$,
alors $x$ n'est pas libre dans $!A$ ; par \textbf{Ax3},
$!A\lif\lforall[x][!A]$ est un axiome, et le modus ponens
donne $\Gamma\Proves\lforall[x][!A]$.

Supposons maintenant que $!A$ s'obtienne par modus ponens à
partir de $\Gamma\Proves !B$ et de
$\Gamma\Proves !B\lif !A$. Par hypothèse de récurrence,
$\Gamma\Proves\lforall[x][!B]$ et
$\Gamma\Proves\lforall[x][(!B\lif !A)]$. Mais
\textbf{Ax2} donne aussi
\[
\Gamma\Proves
\lforall[x][(!B\lif !A)]
\lif(\lforall[x][!B]\lif\lforall[x][!A]).
\]
Deux applications du modus ponens donnent donc
$\Gamma\Proves\lforall[x][!A]$, comme souhaité.
\end{proof}

\begin{thm}\ollabel{thm:WeakGen}
(\emph{Généralisation faible sur les constantes})
Si $\Gamma\Proves !A$ et si $c$ est !!a{constant} qui
ne figure pas dans $\Gamma$, alors il existe une variable
$x$ ne figurant pas dans $!A$ telle que
$\Gamma\Proves\lforall[x][\Subst{!A}{x}{c}]$ ; la
démonstration ne fait pas intervenir $c$.
\end{thm}

\begin{proof}
Soit $!A_1,\ldots,!A_n$ une démonstration de $!A$ à partir
de $\Gamma$, avec $!A=!A_n$. Choisissons une variable $x$
ne figurant dans aucune des formules $!A_1,\ldots,!A_n$,
et considérons la nouvelle suite
$\Subst{!A_1}{x}{c},\ldots,\Subst{!A_n}{x}{c}$.
Elle démontre $\Subst{!A}{x}{c}$ à partir de $\Gamma$.
En effet, pour chaque $i=1,\ldots,n$ :
\begin{itemize}
\item si $!A_i$ est un axiome, alors
$\Subst{!A_i}{x}{c}$ en est un également ;
\item si $!A_i\in\Gamma$, alors $c$ n'y figure pas, donc
$\Subst{!A_i}{x}{c}=!A_i\in\Gamma$ ;
\item si $!A_i$ s'obtient par modus ponens à partir de
$!A_j$ et de $!A_j\lif !A_i$, alors
$\Subst{!A_i}{x}{c}$ s'obtient par modus ponens à partir de
$\Subst{!A_j}{x}{c}$ et de
$\Subst{(!A_j\lif !A_i)}{x}{c}
=\Subst{!A_j}{x}{c}\lif\Subst{!A_i}{x}{c}$.
\end{itemize}
La !!{constant} $c$ ne figure plus dans cette nouvelle
suite. Soit $\Gamma'$ l'ensemble des !!{formula}s de
$\Gamma$ qui apparaissent dans la !!{derivation} de
$\Subst{!A}{x}{c}$. La variable $x$ n'est libre dans aucune
!!{formula} de $\Gamma'$, et
$\Gamma'\Proves\Subst{!A}{x}{c}$. La généralisation
(\olref{thm:Generalization}) donne
$\Gamma'\Proves\lforall[x][\Subst{!A}{x}{c}]$ ; la
monotonie permet donc de conclure
$\Gamma\Proves\lforall[x][\Subst{!A}{x}{c}]$.
\end{proof}

\begin{explain}
Nous voulons remplacer, dans le \olref{thm:WeakGen}, l'exigence
que $x$ ne figure pas du tout dans $!A$ par deux conditions
plus faibles : $x$ n'est pas libre dans $!A$ et est
!!{free for} $c$ dans $!A$. Un changement de variable liée
permet d'y parvenir ; nous allons donc d'abord démontrer
ce résultat.
\end{explain}

\begin{lem}\ollabel{lem:free-for}
Si $x$ est libre pour $c$ dans $!A$ et si $y$ n'est pas
libre dans $!A$, alors $x$ est !!{free for} $y$ dans
$\Subst{!A}{y}{c}$.
\end{lem}

\begin{proof}
Si $x$ n'était pas !!{free for} $y$ dans
$\Subst{!A}{y}{c}$, une occurrence libre de $y$ dans
$\Subst{!A}{y}{c}$ se trouverait dans la portée d'un
quantificateur liant $x$. Or, par hypothèse, $y$ n'est
pas libre dans $!A$ : toutes ces occurrences proviennent
donc de la substitution $\Subst{}{y}{c}$. Une occurrence
de $c$ dans $!A$ se trouve alors dans la portée de ce
quantificateur, ce qui contredit le fait que $x$ est
!!{free for} $c$ dans $!A$.\footnote{Correction de l'édition
française : la première substitution de la preuve anglaise
remplace à tort $x$ au lieu de $c$. Le même opérande $c$
est rétabli dans toute la preuve.}
\end{proof}

\begin{lem}[Changement de variable liée]
\ollabel{lem:ChangeBdVar}
Si $x$ et $y$ ne sont pas libres dans $!A$ et sont tous deux
!!{free for} $c$ dans $!A$, alors
$\Proves\lforall[x][\Subst{!A}{x}{c}]
\equiv\lforall[y][\Subst{!A}{y}{c}]$ ; la démonstration
ne fait pas intervenir $c$.
\end{lem}

\begin{proof}
Par hypothèse, $y$ est !!{free for} $c$ dans $!A$ et
$x$ n'est pas libre dans $!A$. Le lemme précédent,
\olref{lem:free-for}, montre donc que $y$ est !!{free for}
$x$ dans $\Subst{!A}{x}{c}$. Par \textbf{Ax1},
\[
\lforall[x][\Subst{!A}{x}{c}]
\lif\Subst{!A}{x}{c}\Subst{}{y}{x}
\]
est un axiome.\footnote{Correction de l'édition française :
la source place l'implication entière sous le quantificateur
universel. L'axiome d'instanciation quantifie seulement son
antécédent, comme ici.}
Puisque $x$ n'est pas libre dans $!A$,
$\Subst{!A}{x}{c}\Subst{}{y}{x}=\Subst{!A}{y}{c}$.
Ainsi,
\[
\Proves\lforall[x][\Subst{!A}{x}{c}]
\lif\Subst{!A}{y}{c}.
\]
En ajoutant l'antécédent comme hypothèse et en appliquant le
modus ponens, on obtient
$\lforall[x][\Subst{!A}{x}{c}]\Proves\Subst{!A}{y}{c}$.
Comme $y$ n'est pas libre dans $!A$, il n'est pas non plus
libre dans $\lforall[x][\Subst{!A}{x}{c}]$.
La généralisation donne
$\lforall[x][\Subst{!A}{x}{c}]
\Proves\lforall[y][\Subst{!A}{y}{c}]$,
puis le théorème de la déduction donne
$\Proves\lforall[x][\Subst{!A}{x}{c}]
\lif\lforall[y][\Subst{!A}{y}{c}]$.
La démonstration de l'implication inverse,
$\Proves\lforall[y][\Subst{!A}{y}{c}]
\lif\lforall[x][\Subst{!A}{x}{c}]$,
est symétrique ; la règle Taut permet donc de conclure.
\end{proof}

\begin{thm}[Généralisation forte sur les constantes]
\ollabel{thm:StrongGen}
Si $\Gamma\Proves !A$, si la !!{constant} $c$ ne figure pas
dans $\Gamma$ et si $x$ n'est pas libre dans $!A$, mais est
!!{free for} $c$ dans $!A$, alors
$\Gamma\Proves\lforall[x][\Subst{!A}{x}{c}]$.
\end{thm}

\begin{proof}
Puisque $\Gamma\Proves !A$ et que $c$ ne figure pas dans
$\Gamma$, la généralisation faible donne !!a{variable} $y$
ne figurant pas dans $!A$ telle que
$\Gamma\Proves\lforall[y][\Subst{!A}{y}{c}]$.
La variable $y$ n'est donc pas libre dans $!A$ et est aussi
!!{free for} $c$ dans $!A$. Par hypothèse, $x$ satisfait ces
mêmes conditions. Le lemme de changement de variable liée
s'applique, et donne
$\Proves\lforall[x][\Subst{!A}{x}{c}]
\equiv\lforall[y][\Subst{!A}{y}{c}]$.
On en déduit $\Gamma\Proves\lforall[x][\Subst{!A}{x}{c}]$.
\end{proof}

\begin{thm}\ollabel{thm:provability-quantifiers}
On suppose que $t$ est libre pour $x$ dans $!A(x)$.\footnote{
Précision de l'édition française : la source omet cette
condition, qui empêche la capture de variables lors des
instanciations ci-dessous.}
\begin{tagenumerate}{prvEx,prvAll}
\tagitem{prvEx}{Si $\Gamma\Proves !A(t)$, alors
$\Gamma\Proves\lexists[x][!A(x)]$.}{}
\tagitem{prvAll}{Si $\Gamma\Proves\lforall[x][!A(x)]$,
alors $\Gamma\Proves !A(t)$.}{}
\end{tagenumerate}
\end{thm}

\begin{proof}
\begin{tagenumerate}{prvEx,prvAll}
% La source laisse le cas existentiel en exercice.
\tagitem{prvEx}{Exercice.}{}
\tagitem{prvAll}{Supposons que
$\Gamma_0\Proves\lforall[x][!A(x)]$.

\begin{derivation}
1. & $\Gamma_0\Proves\lforall[x][!A(x)]$ & hyp. \\
2. & $\Gamma_0\Proves\lforall[x][!A(x)]\lif\Subst{!A}{t}{x}$ & Ax1 \\
3. & $\Gamma_0\Proves\Subst{!A}{t}{x}$ & MP 1, 2 \\
\end{derivation}
La monotonie permet de conclure pour~$\Gamma$.}{}
\end{tagenumerate}
\end{proof}












\end{document}
